\chapter{Finite Gauge Theory And State-Sum TQFT}
\label{ch:tqft-finite-gauge-state-sum}

\ConventionsEuclidean

Finite gauge theory is the cleanest laboratory for the relation between
path-integral sums, extended functoriality, defects, and anomalies.  The
fields form finite groupoids, the sums are literal finite sums, and the
gluing identities can be checked without analytic regularization.  This
chapter develops the untwisted and Dijkgraaf--Witten twisted theories as
TQFTs.

\section{Fields As A Groupoid}

Let \(G\) be a finite group and let \(M\) be a compact \(D\)-manifold.  The
field groupoid is
\[
  \operatorname{Bun}_G(M),
\]
whose objects are principal \(G\)-bundles \(P\to M\) and whose morphisms are
bundle isomorphisms.  The untwisted partition function is the groupoid
cardinality
\[
  Z_G(M)
  =
  \sum_{[P]\in\pi_0\operatorname{Bun}_G(M)}
  \frac{1}{|\operatorname{Aut}(P)|}.
\]
If \(M\) is connected and \(P\) is represented by a homomorphism
\(\rho:\pi_1(M)\to G\), then
\[
  \operatorname{Aut}(P)\cong C_G(\rho),
\]
the centralizer of the image of \(\rho\).  Hence
\[
  Z_G(M)
  =
  \frac{|\operatorname{Hom}(\pi_1(M),G)|}{|G|}
\]
for connected \(M\).

\section{Twists By Cocycles}

Let
\[
  \omega\in Z^D(BG,U(1))
\]
be a normalized group cocycle.  A bundle \(P\to M\) is classified by a map
\[
  f_P:M\to BG
\]
up to homotopy.  The twisted partition function is
\[
  Z_{G,\omega}(M)
  =
  \sum_{[P]}
  \frac{1}{|\operatorname{Aut}(P)|}
  \exp\left(2\pi i\int_M f_P^*\omega\right).
\]
Changing \(\omega\) by a coboundary changes the integrand by a boundary term;
on closed \(M\) the partition function depends only on the cohomology class
\[
  [\omega]\in H^D(BG,U(1)).
\]
On manifolds with boundary, a choice of trivialization or boundary state is
part of the datum.

\section{Triangulated State Sum}

Choose a triangulation \(K\) of \(M\).  A flat lattice gauge field assigns
group elements \(g_{ij}\in G\) to oriented edges satisfying
\[
  g_{ij}g_{jk}=g_{ik}
\]
on each oriented triangle.  Gauge transformations assign \(h_i\in G\) to
vertices and act by
\[
  g_{ij}\mapsto h_i g_{ij}h_j^{-1}.
\]
For a cocycle \(\omega\), the state-sum weight is a product over oriented
\(D\)-simplices:
\[
  W(g)=\prod_{\sigma^D}
  \omega(g_{\sigma})^{\operatorname{sgn}(\sigma)} .
\]
The cocycle condition \(\delta\omega=1\) is precisely the algebraic identity
needed for invariance under Pachner moves.  Thus triangulation independence
is a finite calculation in group cohomology.

\section{State Spaces}

For a closed \((D-1)\)-manifold \(\Sigma\), the Hilbert space is the vector
space of functions on the boundary field groupoid twisted by the boundary
line:
\[
  \mathcal H(\Sigma)
  =
  \Gamma\bigl(\operatorname{Bun}_G(\Sigma),\mathcal L_\omega(\Sigma)\bigr).
\]
In the untwisted theory this is the space of class functions on
\(\operatorname{Bun}_G(\Sigma)\).  For a bordism \(M:\Sigma_0\to\Sigma_1\), the linear map
\[
  Z(M):\mathcal H(\Sigma_0)\to\mathcal H(\Sigma_1)
\]
is obtained by summing over bulk bundles restricting to the boundary bundles,
with automorphism factors.  Gluing follows from Fubini's theorem for finite
groupoid cardinality.

\section{Extended Objects}

Codimension-one defects between \(G\)- and \(H\)-gauge theories are described
by finite correspondences of groupoids.  A simple class is a subgroup
\[
  K\subset G\times H
\]
with possible cocycle decoration.  Boundary conditions are module categories
over the fusion category associated to \(G\) and \(\omega\).  In three
dimensions, line operators form the Drinfeld center
\[
  Z(\mathrm{Vect}^{\omega}_G),
\]
the category of representations of the twisted quantum double.  Fusion and
braiding are finite tensor-category operations.

\section{Anomaly Boundary}

A \(D\)-dimensional anomalous finite symmetry theory can be realized as a
boundary of a \((D+1)\)-dimensional finite gauge theory whose cocycle is the
anomaly.  If
\[
  \alpha\in Z^{D+1}(BG,U(1)),
\]
then a boundary theory with \(G\) symmetry has partition functions valued in
the transgressed line over \(\operatorname{Bun}_G(M)\).  Gauging is possible only after this
line is trivialized or after a bulk is kept as part of the system.

\begin{openproblem}[Finite models]
\label{op:extended-finite-models}
Construct extended finite-gauge TQFTs with boundaries, defects, junctions,
twists, and anomaly inflow as functors on a chosen bordism
\((\infty,D)\)-category, including comparison with the state-sum formulas
above.
\end{openproblem}
